% !TeX root = ../navier-stokes-manuscript.tex
% ============================================================
% 1.1 Clay Requirements and Definitions
% ============================================================
\subsection{Clay Requirements and Definitions}\label{sub:ClayRequirements}

Let \(\Omega\) be either \(\mathbb R^3\) or the periodic domain \(\mathbb T^3=\mathbb R^3/\mathbb Z^3\).
Fix a viscosity \(\nu>0\), set the external force to zero, and choose a smooth divergence-free initial velocity \(u_0\).
On \(\mathbb R^3\), require \(u_0\) and all of its spatial derivatives to decay rapidly at infinity; on \(\mathbb T^3\), require periodicity instead.
The velocity \(u(x,t)\in\mathbb R^3\) and pressure \(p(x,t)\in\mathbb R\) must satisfy
\[
  \partial_tu+(u\cdot\nabla)u+\nabla p
  =
  \nu\Delta u,
  \qquad
  \nabla\cdot u=0,
  \qquad
  u(x,0)=u_0(x).
\]
The local theory produces a unique smooth solution on a short interval, and \(T_*\in(0,\infty]\) denotes the largest time to which that smooth solution can be continued.
The positive existence-and-smoothness problem studied in this paper is exactly whether
\[
  T_*=\infty
\]
for every admissible \(u_0\).

Each term in the equation belongs to the same fluid motion.
The derivative \(\partial_tu\) records the change of velocity at a fixed spatial point, while \((u\cdot\nabla)u\) records the change caused by carrying velocity through a nonuniform field.
The pressure gradient transmits the constraint imposed by incompressibility, and \(\nu\Delta u\) transfers momentum across spatial variation.
The condition \(\nabla\cdot u=0\) says that the flow preserves volume.
The problem asks whether these coupled effects can evolve the initial fluid forever without producing a loss of smoothness.

\divider{Smoothness}

\begin{definition}[Smoothness]\label{def:Smoothness}
A velocity--pressure pair \((u,p)\) is smooth through a finite time \(T>0\) when
\[
  u\in C^\infty\!\left(\Omega\times[0,T];\mathbb R^3\right),
  \qquad
  p\in C^\infty\!\left(\Omega\times[0,T]\right).
\]
Equivalently, every mixed spatial and temporal derivative of every finite order exists and is continuous through \(t=T\).
The pair is smooth for all time when it is smooth through every finite \(T>0\).
\end{definition}

The word smooth has a precise all-orders meaning.
It does not require one numerical constant to bound every derivative order simultaneously.
It requires each finite derivative order to remain well defined and continuous whenever that order is examined.

\divider{Regularity}

\begin{definition}[Regularity]\label{def:Regularity}
Regularity describes the type and degree of control available for a solution.
In the classical scale, \(C^0\) means continuity, \(C^1\) adds continuous first derivatives, and \(C^k\) requires all mixed derivatives through total order \(k\).
Smoothness is the all-orders case \(C^\infty\).
Sobolev and mixed space--time bounds are other forms of regularity, expressed through norms that measure integral control of the field and its derivatives.
\end{definition}

A regularity statement must name its level.
Continuity of the velocity does not by itself control its gradient, and control of finitely many derivatives does not by itself give \(C^\infty\).
The historical continuation criteria below matter because they identify levels of regularity strong enough to recover all remaining finite orders from the equations.

\divider{Continuation and Maximal Time}

\begin{definition}[Continuation]\label{def:Continuation}
Suppose \((u,p)\) solves the Navier--Stokes equations on \(\Omega\times[0,T)\).
A continuation past \(T\) is a velocity--pressure pair \((\widetilde u,\widetilde p)\) solving the same equations with the same viscosity on \(\Omega\times[0,T+\varepsilon)\) for some \(\varepsilon>0\), agreeing with \((u,p)\) before \(T\), and belonging to the required regularity class on the larger interval.
\end{definition}

For the Clay problem, the required class is \(C^\infty\).
The adjective maximal means that the smooth solution has no such continuation beyond \(T_*\).
Thus \(T_*<\infty\) can occur only if the solution loses enough regularity near \(T_*\) to defeat every applicable local restart.

\divider{Finite Energy}

On the whole space, the required solution also has bounded kinetic energy:
\[
  \sup_{t\geq0}
  \int_{\mathbb R^3}|u(x,t)|^2\,dx
  <
  \infty.
\]
This condition keeps the total motion physically finite, yet it is much weaker than smoothness because it controls the velocity in \(L^2\) without directly controlling its pointwise values or higher derivatives.
The distinction between finite energy and smoothness is where the historical estimates begin.

\divider{What a Loss of Smoothness Means}

At a finite endpoint \(T\), \(C^k\) regularity is lost when at least one required mixed derivative through order \(k\) fails to extend continuously to \(\Omega\times\{T\}\).
That derivative may become unbounded, or it may remain bounded while approaching incompatible limiting values.
A failure at any finite order prevents continuation in \(C^\infty\).
The formal problem is whether the maximal smooth time can ever be finite.
